% Kaiyu Yang — Resume (2 pages)
% Build: latexmk -pdf Kaiyu_Yang_Resume.tex
\documentclass[10.5pt]{article}

\usepackage[margin=0.75in]{geometry}
\usepackage[T1]{fontenc}
\usepackage[utf8]{inputenc}
\usepackage{XCharter}            % matches the website's Charter
\usepackage[svgnames]{xcolor}
\usepackage{enumitem}
\usepackage{titlesec}
\usepackage[hidelinks]{hyperref}
\usepackage{microtype}

\definecolor{accent}{HTML}{175F53}
\hypersetup{colorlinks=true, urlcolor=accent, linkcolor=accent}

\pagestyle{empty}
\setlength{\parindent}{0pt}
\setlength{\parskip}{0pt}

\titleformat{\section}
  {\normalsize\bfseries\scshape\color{accent}}
  {}{0pt}{}
  [\vspace{-0.55em}{\color{accent!40}\rule{\textwidth}{0.4pt}}]
\titlespacing*{\section}{0pt}{1.1em}{0.55em}

\newcommand{\entry}[4]{% org, location, title, dates
  \textbf{#1}, #2 \hfill #4\\
  \textit{#3}\par}

\newcommand{\pub}[3]{% title+link, authors, venue
  \item \href{#3}{\textbf{#1}}\\ {\small #2}}

\newcommand{\pubnolink}[2]{%
  \item \textbf{#1}\\ {\small #2}}

\setlist[itemize]{leftmargin=1.15em, itemsep=0.18em, parsep=0pt, topsep=0.25em}

\begin{document}

{\LARGE\bfseries Kaiyu Yang}\hfill
\begin{tabular}[b]{r@{}}
  \small kaiyuy@alumni.princeton.edu\\[-1pt]
  \small \href{https://yangky11.github.io}{yangky11.github.io} \, · \,
  \href{https://scholar.google.com/citations?user=FciCu4EAAAAJ}{Google Scholar} \, · \,
  \href{https://github.com/yangky11}{github.com/yangky11}
\end{tabular}

\vspace{0.35em}
{\color{accent!40}\rule{\textwidth}{0.8pt}}
\vspace{0.35em}

Research leader working on capable, verifiable AI: the mathematical capabilities of large
language models across the full training stack, and verification for AI systems and agents.
Working on machine learning for formal mathematics and theorem proving since 2018,
from early benchmarks (CoqGym) to widely used open-source platforms (LeanDojo) and
open state-of-the-art provers (Goedel-Prover).

\section{Experience}

\entry{Apodex}{Redwood City, CA}{Lead Scientist}{Mar 2026 -- Present}
\begin{itemize}
  \item Lead a research team working on (1) the mathematical capabilities of LLMs across the
        full stack --- pre-training, mid-training, post-training, and agents --- and (2)
        verification: formal verification for LLM training infrastructure, and verifiers that
        make agents efficient and reliable on long-horizon tasks.
  \item Set the lab's research agenda and direct external academic collaborations on
        theorem proving and verified code generation.
\end{itemize}

\vspace{0.5em}
\entry{Meta Fundamental AI Research (FAIR)}{New York, NY}{Research Scientist}{Jun 2024 -- Dec 2025}
\begin{itemize}
  \item Co-authored \textit{Formal Mathematical Reasoning: A New Frontier in AI}
        (ICML 2025 position paper, spotlight; also in \textit{Communications of the ACM}),
        which lays out the research agenda connecting LLMs with formal methods.
  \item Advised the Goedel-Prover series (COLM 2025, ICLR 2026) --- open-source theorem
        provers that set the state of the art among open models --- and ProofOptimizer
        (ICLR 2026), a model that simplifies machine-generated proofs without human demonstrations.
  \item Collaborations on verified code generation (Verina, ICLR 2026) and neuro-symbolic
        olympiad problem solving (ICLR 2025, CAV 2025).
\end{itemize}

\vspace{0.5em}
\entry{California Institute of Technology}{Pasadena, CA}{Computing, Data, and Society Postdoctoral Fellow}{Sep 2022 -- May 2024}
\begin{itemize}
  \item Built \href{https://www.leandojo.org/}{LeanDojo} (NeurIPS 2023, oral), an open-source
        platform for LLM-based theorem proving in Lean, and Lean Copilot, which runs LLM
        inference natively inside the Lean proof assistant; both are in regular use in the field.
  \item Co-organized the MATH-AI workshop and the tutorial on machine learning for theorem
        proving at NeurIPS 2023.
\end{itemize}

\section{Education}

\entry{Princeton University}{Princeton, NJ}{PhD in Computer Science, advised by Jia Deng}{2022}
\vspace{0.35em}
\entry{University of Michigan}{Ann Arbor, MI}{MS in Computer Science and Engineering}{2018}
\vspace{0.35em}
\entry{Tsinghua University}{Beijing, China}{B.Eng.\ in Computer Science and B.S.\ in Mathematics}{2016}

\section{Selected Publications}

Full list on \href{https://scholar.google.com/citations?user=FciCu4EAAAAJ}{Google Scholar}
(30+ papers; NeurIPS, ICML, ICLR, COLM, CACM, EMNLP, CVPR, ICCV, ECCV, CAV).

\begin{itemize}
  \pub{LeanDojo: Theorem Proving with Retrieval-Augmented Language Models}
      {K.~Yang, A.~Swope, A.~Gu, R.~Chalamala, P.~Song, S.~Yu, S.~Godil, R.~Prenger, A.~Anandkumar. NeurIPS 2023 (oral).}
      {https://www.leandojo.org/}
  \pub{Formal Mathematical Reasoning: A New Frontier in AI}
      {K.~Yang, G.~Poesia, J.~He, W.~Li, K.~Lauter, S.~Chaudhuri, D.~Song. ICML 2025 (spotlight); Communications of the ACM.}
      {https://arxiv.org/abs/2412.16075}
  \pub{Goedel-Prover: A Frontier Model for Open-Source Automated Theorem Proving}
      {Y.~Lin*, S.~Tang*, B.~Lyu, J.~Wu, H.~Lin, K.~Yang, J.~Li, M.~Xia, D.~Chen, S.~Arora, C.~Jin. COLM 2025.}
      {https://blog.goedel-prover.com/}
  \pub{Goedel-Prover-V2: Scaling Formal Theorem Proving with Scaffolded Data Synthesis and Self-Correction}
      {Y.~Lin*, S.~Tang*, B.~Lyu*, Z.~Yang*, et al. ICLR 2026.}
      {https://blog.goedel-prover.com/}
  \pub{Verina: Benchmarking Verifiable Code Generation}
      {Z.~Ye, Z.~Yan, T.~Kasriel, J.~He, K.~Yang, D.~Song. ICLR 2026.}
      {https://verina.io/}
  \pub{Lean Copilot: Large Language Models as Copilots for Theorem Proving in Lean}
      {P.~Song, K.~Yang, A.~Anandkumar. NeuS 2025.}
      {https://arxiv.org/abs/2404.12534}
  \pub{Proving Olympiad Inequalities by Synergizing LLMs and Symbolic Reasoning}
      {Z.~Li*, Z.~Li*, W.~Tang, X.~Zhang, Y.~Yao, X.~Si, F.~Yang, K.~Yang\textsuperscript{\dag}, X.~Ma\textsuperscript{\dag}. ICLR 2025.}
      {https://arxiv.org/abs/2502.13834}
  \pub{Generating Natural Language Proofs with Verifier-Guided Search}
      {K.~Yang, J.~Deng, D.~Chen. EMNLP 2022 (oral).}
      {https://arxiv.org/abs/2205.12443}
  \pub{Learning to Prove Theorems via Interacting with Proof Assistants}
      {K.~Yang, J.~Deng. ICML 2019.}
      {https://arxiv.org/abs/1905.09381}
\end{itemize}

\section{Open Source}

Core developer of \href{https://github.com/lean-dojo/LeanDojo}{LeanDojo},
\href{https://github.com/lean-dojo/ReProver}{ReProver},
\href{https://github.com/lean-dojo/LeanCopilot}{Lean Copilot}, and
\href{https://github.com/princeton-vl/CoqGym}{CoqGym} --- open-source tooling for
machine learning with the Lean and Coq proof assistants.

\section{Community Leadership}

\begin{itemize}
  \item Co-organizer, Workshop on Mathematical Reasoning and AI (MATH-AI) at NeurIPS 2023,
        2025, and 2026; co-organizer of the NeurIPS 2023 tutorial on machine learning for
        theorem proving.
  \item Area Chair: ICML 2025 \& 2026, NeurIPS 2026, ECCV 2024. Reviewer for NSF panels and
        ERC Advanced Grants.
  \item Guest co-instructor, \textit{Advanced Large Language Model Agents} (UC Berkeley \& MOOC,
        2025); guest lecturer at Caltech and CUHK.
  \item Invited talks at OpenAI, Google, Meta FAIR, the Simons Institute, UC Berkeley, Stanford,
        CMU, and other universities.
  \item Mentored 8+ students and junior researchers, now in PhD programs (Toronto, Cornell,
        CUHK, UIUC) and at industry labs.
\end{itemize}

\section{Recognition}

\begin{itemize}
  \item Siebel Scholar (2022) --- 42 computer science graduate students awarded annually worldwide.
  \item Oral presentations at NeurIPS 2023 and EMNLP 2022; ICML 2025 spotlight; NeurIPS 2020 spotlight.
  \item Outstanding Reviewer, CVPR 2020 \& 2021 (top 20\%).
  \item Research covered by Scientific American and Wired.
\end{itemize}

\vspace{0.4em}
{\small * equal contribution \quad \textsuperscript{\dag} equal advising}

\end{document}
